\section{Death and Deliverance}

\begin{quotex}
This is second and concluding installment of the chapter ``La Morte'' from Dio e il Poeta by \textbf{Guido De Giorgio}. Here we see De Giorgio moving past salvation to liberation (or deliverance as Guenon translates it). Evola draws upon similar traditional teachings in his distinction between the path of the ancestors and the path of the gods. For liberation, even Heaven is a prison. De Giorgio is forced to use apophatic language\footnote{\url{http://orthodoxwiki.org/Apophatic_theology}} to describe this process of deliverance. He leads us to the God beyond God, non-Being beyond Being, or in Vedantic terms, Brahman behind Ishwara.
\end{quotex}


Divine knowledge is like a voice that passes by and does not return: whoever does not subscribe to it in freedom, will have his own negation as fate, i.e., the darkness, the void, death. And we arrive in this way at the crucial point: what is Death? Is it truly a cessation \emph{in toto}, an irremediable end, a no longer existing, or a simple, even if complex, break in the natural process of transformation? For common belief, it is the end, the stop, which follows what is not known nor does it want to be fixated on it, but that is, in sum, a not-life, a no longer being what one was. This opinion is common as much to the deniers of God as well as to the lukewarm adherents of Religion. But for Religion, instead, Death is only a crucial point from which another life begins, one more certain, more secure, and truer than that false life that is experienced in the world, the world being a bridge, a necessary passage that conditions the future destiny of man who, we affirm it, is its judge to the extent he \emph{wants} to be it.

How, in fact, could there be light for the man who obstinately rejects the light and how would there be salvation for the man who is convinced that he must be saved from nothing and nobody has to be saved in the inexplicable drama of life? Certainly, when a man passes from the essence of God, from His being all that is, to men's \emph{post mortem} destiny, to the diversity of that destiny and above all to the rather vast exclusion of most from the life in God, he is seized by doubts, he is chased by protests, he is here ontologically out of place, as in the face of the partial negation of a total Affirmation, to an exclusive part of the All, to a tragic void in an absolute fullness, to an inexplicable residue of the divine Universality, to a clashing collision, to a radical injustice. But that depends on the limitations of our vision that is only human, and therefore is in a certain sense
\emph{passionate}, and of our reason that is confined to fixed fields of

logic.

One thinks only that God offers Himself to those who want Him and denies Himself to those who deny him and that, at least and only in this, that ``reciprocity'' is manifestly just, in asking and giving, as a man can claim nothing, if, when he returned home in the evening, he did not find what he didn't want to find, rejecting the unpopular guest to whom he always denied access to his own home, slamming the door in his face.

God cannot give himself to those who do not want Him, to insist with those who are deaf and blind. Religion requires men of good will who adhere to and want to adhere to the Law of God. Whoever rejects it is judged by himself before being judged by God. This personal way of seeing things can be very useful only to those who adopt it, even if necessarily limited and partial. Because, examining the thing from a higher and truer point of view, the problem of his own salvation vanishes in the face of that more universal idea of \emph{liberation} which is the dissolution \emph{in toto} of the human-cosmic illusion in God. It corresponds to what St Paul says; when he speaks of the crucifixion of the world, he alludes precisely to the entire and universal process of liberation that is a definitive escape from every limitation, even the supernatural, for the realization of God in the negation of everything, even of Paradise.

But this integrating God in oneself and for oneself, this centering Him in His purest Essence, which is the exclusion of every perfection, for the perfection of all the perfections, this ontological rarefaction that is the negation of every affirmation, this surpassing of every praise, of every eulogy, of every superlative, is truly the intuition of the True God, of the Deity about Whom nothing can be said, not even that He is, if being is only being, and if beyond being, there is the unthinkable absolute, the unintelligible absolute, the absolute ineffability of the Ineffable. In this total super-eminence, God is more than God, since even God is surpassed in the vertiginous exclusion of every perfection. The realization of so much muchness that coincides with the mininumness of every minimum, it is the highest degree of divine realization and it is the silence of silence in silence: here God says of Himself: \emph{Ego sum quod non est}, the neuter absorbing in his indeterminative infinity every personal veil, every aseity, every formulation is only tautological.

There is no more mirror---\emph{Ego sum qui sum}---not even of God in God, not a name, not a voice, not a song: nothing. The Being that is, is the non-being that is. It suppresses the last, the universal determination, which is that of being, of aseity, nothing remains if not the Nothing, the non-Being, the Ineffable of the Ineffabile. This is liberation and those who attain to it will find ``that even Paradise is a prison'' and that there is no longer anything to attain for those who ask for nothing because they can have nothing more. This is not an absurdity of words, but rather a radical absurdity, a divine absurdity that climbs over every limit, every peak, that fills up every chasm, that sinks in the intimate, hidden beauty of the Infinite. For so high a liberation, so total, Death is nothing, barely a transit, not even a passage, but the door through which the Wise, the Saint, abandons the cosmic prison, climbs through all the heavens, and penetrates into the ultimate, \emph{Coelum Coeli}, where God is more than God, He is the infinite totality of negation. Here all the songs are silent and the Abyss alone is silent.

Death is a great master of wisdom and virtue if it is known like that, to what sublimities and what terrible ruins it leads, if one familiarizes oneself with its advent, if it is obtained from God who prepares us to receive it as the necessary means for the exodus from this ``false world'' and the access point to the highest degrees of Being, to the boundless possibilities of new worlds where the soul can sing its loftiest victories. God's true wisdom prepares us to accommodate it as the liberation from the corporeal burden, as the abandonment of the cadaver that we drag on earth, since the life of the multitude, deprived of God's light, is foolish, dull, a worthless necrobia, a dragging of artificial superstructures that hide vermin, the putrefaction of a dead man who dreams.

Dying before dying prepares us for death, feeling oneself a cadaver before becoming one, separating before detachment, sanctifying the corpse, purifying it, making it the marble temple of the Living God. The sanctification of the body is the sanctification of life and the sanctification of the soul that only the Holy calls Holy, i.e., God in us. The divine integration, the knowledge that God alone thinks, feels, and works in us, is the most certain means to attain freedom and serenity at the transition.

Whatever pain, whatever sacrifice, St Paul says, is nothing in the face of the infinite richness of the prize that is now pure of God in comparison to the meager and false earthly riches. Freeing oneself from the tyranny of the I, one can gradually and laboriously attain to peace, to the triumph of the eternal in us, when all the interior springs flow with fresh water. It is a type of progressive denudation, of spiritual desquamation, of the gradual frustration of transience, of no longer feeling the demon that roars in us, the tormenting hammering of the I that arrogantly grabs the openings of the soul and twists them to pain, to crash, to laugh, to tears. To empty life of every content, to give it back to God, through contemplation, prayer, invocation, pneumatic elevation, radiant of Love, the inward conversation with Saints and Masters who become our Guides, our helpers, our mirrors of virtue. This reducing oneself, this abolishing the inner thunder, this listening to the perennial sources of peace flowing, this flowing of silence, this caressing the void in its fullness, this perceiving the darkness of oneself until the Light bubbles, this darkening, this crumbling and incorporeal re-emerging, this confident abandonment in God, this acceptance of all only for the love of the Lord, this No to the world and this Yes to God, this reflecting oneself in other-than-oneself, no longer recognizing one's own face, and ignoring everything except for Him, here is how rebirth is realized after the end, when Death, death killed, is the triumphal song of the soul that rests in God.

See ⇐ Part 1\footnote{\url{https://gornahoor.net/?p=5975}}


\flrightit{Posted on 2013-02-26 by Aeneas}

\begin{center}* * *\end{center}

\begin{footnotesize}
\begin{sffamily}
\texttt{Izak on 2013-02-27 at 16:21 said:}

This really is very beautiful. I realize that your `musings' posts get more comments, but I think the reason is because these translations are harder to digest and demand a higher degree of contemplation, so everyone risks looking like a fool so much more before saying anything.

But I think everyone appreciates these very invaluable translations, perhaps more than anything else on the site. The fifth paragraph is especially nice. I'm assuming that its poetry is even more strong in the original Italian than it is here.

\hfill

\texttt{Cologero on 2013-02-27 at 22:12 said:}

I appreciate your interest, Izak. As you may imagine, it is not necessarily easy to translate, to find the right word, to be sure that the meaning is correct. It is actually a transformative process to do so, and I hope the readers of the translation will experience the same process. It is important to avoid being presumptuous, i.e., by assuming the meaning is obvious and clear cut. Some facebook comments claim the text is ``life denying''; if so, their lives must be quite spectacular.

\hfill

\texttt{george on 2013-02-28 at 00:22 said:}

The perspective is so thoroughly antimodern that it is hard to contemplate without letting go of modern perspectives.

\hfill

\texttt{Izak on 2013-02-28 at 14:49 said:}

Yeah, I'd say you did about the best translation that the English language can permit, at least for that paragraph. I can get a sense of what Guido was doing with language.

It's really nice to see Guido de Giorgio in English either way. I can see why people would call this passage ``life denying,'' but real metaphysics always kind of does walk the tightrope, doesn't it! It's very hard to say the truth without appearing Manichaean or pessimistic or whatever. Not to totally switch topics, but I always did appreciate when John Woodroffe took pains to explain that ``maya'' does not mean ``illusion'' in the way Westerners think about it, with negative or moralistic connotations. If only more people could be alleviate of such confusions!

\hfill

\texttt{Cologero on 2013-03-03 at 08:00 said:}

Someone called me last week and in the course of the conversation, he claimed that ``God would not send him to hell.'' At the time, I thought it was an odd comment from someone who frequents Gornahoor, not least of all because of the rather blasphemous view of God as some sort of remote ``sky god'', occasionally poking around in the world rather arbitrarily and even unjustly.

Coincidentally, since this particular post was scheduled, rather than attempt to answer him directly, I asked him to read this post when it appeared to see if it addressed his particular concern. I assume it has or else he lost interest. De Giorgio claims that in the post mortem state, the soul will itself seek its own level based on how much ``light'' it can bear. He claims to be repeating traditional teachings.

It is true that Plotinus said that different souls would exist on different levels, reflecting their large diversity, Dante mapped out these levels in the Divine Comedy. It seems obvious that a soul attached to a certain vice will have an affinity to a certain level of hell. It is pointless to put him on a higher level since he will just fall back to his own lever; there is no need to ``send'' him there as though it were against his own will. This is not a theory, it can be verified in your own consciousness, at least for those who have achieved the proper level of self-awareness.

To avoid such a fate, the important issue is to understand the ``light''. It is well worth the effort. According to traditional teachings, the Self is self-effulgent.

\hfill

\texttt{Michael on 2020-03-24 at 10:14 said:}

Just the way GdG forms this writing is amazing, even more so is his light which the writing attempts to convey.

And yet he gives us the parts of a process, specifically the last two paragraphs. All of the statements made there convey exercises to be done by us the reader. Perhaps more can be elaborated on them, will sit with them to digest.

\end{sffamily}
\end{footnotesize}

